\chapter{Contexto científico y objetivos del proyecto}
\label{chap:contexto_objetivos}
\section{Necesidad de una estimación espacial}
En intervencionismo, la dosis cutánea depende de la secuencia de proyecciones, la posición del paciente y la coincidencia de campos. Dos procedimientos con igual kerma acumulado pueden concentrar la irradiación en una zona o repartirla entre regiones distintas. La distribución espacial aporta por tanto información que no contiene un único indicador integrado \citep{aapm2021,malchair2020}.

La motivación clínica del mapeo no convierte automáticamente el mapa en un predictor de lesión. La respuesta tisular depende de dosis, distribución, tiempo y factores individuales; DermaDose no incorpora un modelo radiobiológico. Su problema actual es dosimétrico y geométrico. La presentación de resultados debe evitar traducir una escala de colores a un diagnóstico.


\section{Familias metodológicas}
Pueden distinguirse modelos basados en traslado del kerma y correcciones, reconstrucciones anatómicas con transporte de partículas y medidas directas de distribución. La primera familia reduce el coste y facilita interpretar cada factor; la segunda puede incorporar heterogeneidades y dispersión espacial; la tercera depende de la respuesta, calibración y colocación de los detectores. Estas familias pueden complementarse en una estrategia de validación.

DermaDose pertenece actualmente a la primera. Su malla es una frontera geométrica donde se evalúa la entrada. No es una discretización volumétrica de tejidos. El factor de retrodispersión resume una respuesta calculada en agua; no genera trayectorias en el cuerpo representado.




\section{Justificación y pregunta del trabajo}

En los procedimientos intervencionistas guiados por fluoroscopia, una dosis absorbida local en piel del orden de 2~Gy puede dar lugar a reacciones tisulares tempranas, como el eritema transitorio. Cuando varias proyecciones irradian una misma región cutánea, sus contribuciones se acumulan y pueden generar un máximo localizado. Sin embargo, los indicadores de exposición disponibles en el procedimiento, como el kerma en aire acumulado en el punto de referencia, no describen por sí solos la distribución espacial de la dosis sobre la piel ni permiten determinar la localización y magnitud de ese máximo.

Esta limitación motiva el desarrollo de DermaDose, cuyo objetivo es reconstruir la distribución de dosis sobre la superficie cutánea para estimar la dosis pico en piel (PSD) y localizar las regiones más expuestas. La pregunta del trabajo es cómo reconstruir, de manera trazable y reproducible, la contribución de cada evento de irradiación y su acumulación sobre una superficie anatómica, manteniendo explícitas las hipótesis de geometría y respuesta dosimétrica.

El caso de aplicación es el angiógrafo Philips Azurion del Hospital Universitario Marqués de Valdecilla, equipado con un tubo MRC 200 0407.
\section{Objetivo general}
Desarrollar y verificar un calculador de dosis superficial independiente del fabricante que utilice informes RDSR, una superficie tridimensional y una descripción complementaria del equipo, y adaptarlo documental y computacionalmente al Philips Azurion con MRC 200 0407 para sustentar el estudio de procedimientos intervencionistas del hospital.

La independencia se refiere a la implementación y a la trazabilidad del cálculo. No elimina la necesidad de conocer y caracterizar el equipo ni permite asumir compatibilidad automática con cualquier sistema.

\section{Objetivos específicos y resultados evaluables}
\begin{enumerate}[leftmargin=*]
\item Leer y normalizar los eventos RDSR, conservando unidades, identidad y procedencia; diferenciar información ausente, valores inválidos y declaración explícita de ausencia de filtro.
\item Reconstruir la geometría del haz y su primera entrada sobre una malla, incorporar las transformaciones de coordenadas declaradas y acumular las contribuciones sin irradiar artificialmente la superficie de salida.
\item Disponer de superficies de referencia y de un procedimiento de segmentación y exportación desde TC mediante DermaDose Surface, con control de integridad y resolución de la malla.
\item Implementar factores dosimétricos con definiciones físicas compatibles, dominios de interpolación explícitos y trazabilidad bibliográfica. Integrar los datos monoenergéticos de Benmakhlouf con espectros SpekPy para las calidades modeladas.
\item Adaptar la descripción del Azurion al MRC 200 0407: identidad, ángulo anódico, límites del conjunto y separación de filtración fija y añadida. Distinguir datos documentales, parámetros medidos e hipótesis nominales.
\item Producir mapas y métricas de máximo superficial que puedan revisarse y reproducirse, acompañados de la configuración efectiva, las tablas seleccionadas y las limitaciones de interpretación.
\item Verificar propiedades físicas y regresiones mediante pruebas y un caso reproducible de principio a fin. Delimitar estas comprobaciones frente a la validación experimental independiente.
\item Establecer una metodología para analizar datos de procedimientos de Valdecilla: comprobar elegibilidad y calidad de las entradas, describir indicadores y máximos estimados, y estudiar su variabilidad según el procedimiento. La aplicación a una cohorte y las conclusiones estadísticas requieren datos y resultados propios documentados.
\end{enumerate}

